\subsection{Sistemas distribuidos, transparencias ANSA y CAP}
\label{sec:fundamento-belinda}

\subsubsection{Concepto y características de un sistema distribuido}

Un sistema distribuido organiza procesos y recursos en varios equipos conectados
para prestar servicios de manera coordinada. Su diseño requiere decidir qué
componentes se separan, cómo se comunican y qué comportamiento se mantiene
cuando una parte del sistema deja de responder. Van Steen y Tanenbaum explican
que la distribución debe responder a necesidades concretas, como acercar recursos
a los usuarios, compartir servicios o mejorar la continuidad de funcionamiento;
la separación entre equipos también introduce dependencias y fallos parciales
\cite[pp.~3--8]{belinda_vansteen2023_distributed}.

Esta organización presenta varias características relevantes para la práctica.
La comunicación sustituye parte del acceso directo a datos locales por intercambios
entre procesos; las operaciones pueden solaparse y los componentes no siempre
terminan al mismo tiempo. Además, un proceso puede funcionar mientras otro falla
o queda incomunicado. Por ello, conocer el resultado de una operación remota exige
interpretar respuestas, errores y esperas, y proteger el estado compartido frente a
accesos concurrentes \cite[pp.~7--8, 13]{belinda_vansteen2023_distributed}.

Entre los objetivos de diseño se encuentran compartir recursos, ofrecer interfaces
compatibles y facilitar la expansión del servicio. Estos objetivos deben equilibrarse
con el coste de coordinación. Blinowski, Ojdowska y Przybyłek muestran, en su
comparación experimental entre monolitos y microservicios, que distribuir una
aplicación no garantiza una mejora de rendimiento en todas las condiciones. Sus
conclusiones están vinculadas a las aplicaciones y plataformas evaluadas; constituyen
un antecedente para justificar mediciones controladas, no un resultado anticipado
sobre TCP y gRPC \cite[pp.~20371--20372]{belinda_blinowski2022_architecture}.

En esta práctica, los nodos A, B y C son procesos independientes ejecutados en
\texttt{localhost}. Esta configuración reproduce la separación lógica y el intercambio
de mensajes, pero los tres procesos comparten equipo, sistema operativo y recursos
físicos. En consecuencia, permite estudiar mecanismos de comunicación y ordenamiento
de eventos sin representar por completo los retardos, fallos independientes y
condiciones de una red entre máquinas diferentes. Esta delimitación corresponde
al diseño experimental del equipo.

\subsubsection{Comparación con sistemas centralizados y paralelos}

La centralización y el paralelismo describen aspectos distintos de un sistema.
La primera se relaciona con dónde se concentra una función o su control; el segundo,
con la ejecución simultánea de partes de un trabajo. Un servicio lógicamente central
puede tener una implementación física distribuida. A su vez, un clúster puede ser
distribuido y ejecutar un cálculo paralelo mediante paso de mensajes
\cite[pp.~6--7, 33--34]{belinda_vansteen2023_distributed}.
La Tabla~\ref{tab:comparacion-sistemas} sintetiza estas diferencias como criterios
de comparación, sin tratarlos como categorías mutuamente excluyentes.

\begingroup
\small
\renewcommand{\arraystretch}{1.2}
\begin{longtable}{@{}>{\raggedright\arraybackslash}p{2.7cm}>{\raggedright\arraybackslash}p{3.8cm}>{\raggedright\arraybackslash}p{3.8cm}>{\raggedright\arraybackslash}p{4.0cm}@{}}
\caption{Criterios para comparar centralización, paralelismo y distribución. Elaboración propia a partir de \cite[pp.~6--7, 33--34]{belinda_vansteen2023_distributed}.}
\label{tab:comparacion-sistemas}\\
\toprule
Criterio & Centralizado & Paralelo & Distribuido \\
\midrule
\endfirsthead
\toprule
Criterio & Centralizado & Paralelo & Distribuido \\
\midrule
\endhead
Organización & Una función o decisión se concentra en un punto lógico o físico.
& Varias unidades ejecutan partes de un trabajo simultáneamente.
& Procesos y recursos cooperan desde componentes conectados por una red. \\
Comunicación & Puede resolverse mediante llamadas dentro de un proceso o acceso a un servicio central.
& Puede usar memoria compartida o paso de mensajes, según la arquitectura.
& Debe contemplar intercambios entre procesos y las condiciones de la comunicación. \\
Reto principal & Gestionar capacidad, dependencias y continuidad de la función central.
& Dividir tareas y sincronizar resultados con un coste aceptable.
& Coordinar componentes, concurrencia y fallos parciales. \\
Relación con la práctica & Un coordinador de pruebas puede iniciar y supervisar los nodos.
& Los clientes pueden enviar de forma concurrente; eso no garantiza acelerar el trabajo.
& Los tres procesos intercambian mensajes y mantienen estado lógico por nodo. \\
\bottomrule
\end{longtable}
\endgroup

La última fila expresa decisiones del experimento. Distinguir estos criterios evita
confundir el uso de varios procesos con una mejora automática de velocidad o con la
ausencia de puntos centrales de coordinación.

\subsubsection{Las ocho transparencias requeridas por la guía}

La transparencia de distribución busca que usuarios y aplicaciones utilicen un
servicio sin gestionar todos los detalles de la ubicación, representación o
coordinación de sus componentes. Puede apoyarse en una capa de software intermedia
que exponga interfaces comunes \cite[pp.~11--13]{belinda_vansteen2023_distributed}.

Para responder a la clasificación solicitada en la guía, se consideran acceso,
ubicación, migración, reubicación, replicación, concurrencia, fallos y persistencia.
La Figura~1.3 del libro consultado presenta las primeras siete; para persistencia
se utiliza además el análisis de infraestructuras distribuidas de Soiland-Reyes,
Goble y Groth, que incluye esta propiedad en su Tabla~4
\cite[pp.~12--13]{belinda_vansteen2023_distributed}
\cite[tabla~4, pp.~16--17]{belinda_soiland2024_transparency}.
La Tabla~\ref{tab:transparencias} reúne una explicación adaptada de esas categorías.

\begingroup
\small
\renewcommand{\arraystretch}{1.2}
\begin{longtable}{@{}>{\raggedright\arraybackslash}p{3.0cm}>{\raggedright\arraybackslash}p{12.0cm}@{}}
\caption{Transparencias consideradas para el análisis del sistema. Síntesis propia basada en \cite[pp.~11--15]{belinda_vansteen2023_distributed} y \cite[tabla~4]{belinda_soiland2024_transparency}.}
\label{tab:transparencias}\\
\toprule
Transparencia & Aspecto que se busca ocultar \\
\midrule
\endfirsthead
\toprule
Transparencia & Aspecto que se busca ocultar \\
\midrule
\endhead
Acceso & Diferencias de representación y mecanismos de acceso a recursos, mediante interfaces que no obliguen al usuario a manejar cada detalle interno. \\
Ubicación & La localización física del recurso, permitiendo identificarlo de forma lógica sin depender directamente de la máquina que lo aloja. \\
Migración & El cambio de ubicación de un recurso o proceso sin exigir que la aplicación modifique su forma de utilizarlo. \\
Reubicación & El traslado de un recurso mientras está siendo utilizado, procurando conservar las operaciones o interacciones en curso. \\
Replicación & La existencia de varias copias que colaboran para ofrecer un recurso o servicio presentado como una unidad. \\
Concurrencia & Las interferencias entre accesos simultáneos, conservando el comportamiento definido para el recurso compartido. \\
Fallos & La ocurrencia y recuperación de fallos internos, dentro de las condiciones que la implementación realmente pueda manejar. \\
Persistencia & La desactivación y reactivación de recursos, manteniendo para sus usuarios una forma estable de utilizarlos a través de su ciclo de vida. \\
\bottomrule
\end{longtable}
\endgroup

Estas transparencias son objetivos de diseño y deben comprobarse en cada
implementación. Ocultar una operación detrás de una función no elimina los retardos
ni asegura que un fallo pueda recuperarse. El libro advierte que intentar ocultar
por completo la distribución puede aumentar el coste y dificultar que el programador
comprenda los errores \cite[pp.~14--15]{belinda_vansteen2023_distributed}.
En el sistema del equipo, los CSV conservan evidencia experimental, pero ese
almacenamiento no restaura el estado de los nodos al reiniciarlos. La evaluación
de las transparencias implementadas se presenta en el
apartado~\ref{sec:ansa-sistema}.

\subsubsection{Teorema CAP y compromisos de diseño}

CAP expresa un límite para sistemas que mantienen estado compartido a través de
componentes conectados por una red: durante una partición no pueden garantizarse
simultáneamente consistencia fuerte y disponibilidad para todas las solicitudes.
Lee y sus colaboradores retoman este límite y estudian cómo los requisitos de una
aplicación condicionan las decisiones entre consistencia, disponibilidad y latencia
\cite[pp.~138:2--138:3]{belinda_lee2023_cap}.

Para interpretar el teorema deben precisarse sus tres propiedades:
\begin{description}
    \item[Consistencia (C).] Las operaciones observan un estado compatible con una
    única copia lógica y respetan el orden real de las operaciones ya completadas.
    Esta exigencia se relaciona con la linealizabilidad, según la cual una operación
    parece producir su efecto en algún instante entre su inicio y su finalización
    \cite[p.~400]{belinda_vansteen2023_distributed}.
    \item[Disponibilidad (A).] Las solicitudes dirigidas a componentes que no han
    fallado pueden obtener una respuesta válida del servicio. La propiedad no debe
    confundirse con observar un RTT pequeño durante una corrida sin fallos.
    \item[Tolerancia a particiones (P).] El diseño contempla que se interrumpa la
    comunicación entre partes del sistema y debe definir qué servicio conserva
    en esa situación \cite[p.~138:2]{belinda_lee2023_cap}.
\end{description}

Como ejemplo conceptual, supóngase que dos réplicas mantienen una unidad disponible
de un producto. La comunicación entre ellas se interrumpe y una venta confirmada
en la primera reduce su inventario a cero. Si la segunda responde todavía que hay
una unidad, puede atender la consulta, pero no conserva la visión de una única
copia actualizada. Si necesita conocer la actualización antes de responder, deberá
posponer o rechazar ciertas operaciones mientras no pueda comunicarse. El ejemplo
ilustra el compromiso; no corresponde a datos ni a una funcionalidad implementada
en esta práctica. La revisión de Ahmed y colaboradores analiza este problema en
almacenes replicados y discute los costes de coordinación
\cite[pp.~172--173]{belinda_ahmed2023_consistency}.

La prioridad depende del servicio. Lee y colaboradores analizan, por ejemplo, la
coordinación del acceso de vehículos a una intersección: actuar con desacuerdo sobre
su estado puede comprometer la seguridad, por lo que puede ser necesario esperar.
En contraste, explican que la respuesta al pedal de freno no debería quedar bloqueada
por la llegada de un resultado remoto de visión artificial. Son aplicaciones
concretas en las que el coste de esperar y el de actuar con información incompleta
son diferentes \cite[pp.~138:2--138:3]{belinda_lee2023_cap}.

El análisis tampoco termina cuando la red funciona sin particiones. La formulación
PACELC considera entonces el compromiso entre latencia y consistencia: coordinar
réplicas puede añadir intercambios y esperas aun sin una interrupción de red
\cite[p.~172]{belinda_ahmed2023_consistency}.
Por ello, el diseño de una aplicación debe especificar qué datos comparte, qué
operaciones requieren coordinación y qué respuesta permite cuando esa coordinación
no es posible.

En el alcance de esta práctica, TCP y gRPC proporcionarán mecanismos de comunicación.
Elegir uno de ellos no determina por sí solo una clasificación AP o CP: esa decisión
depende de la lógica del servicio y del tratamiento de su estado compartido.
Las pruebas de RTT, concurrencia y relojes de Lamport permitirán estudiar el
comportamiento de los mensajes; sin un protocolo de replicación y un experimento
de particiones definido, no constituirán una demostración experimental de CAP.
